\section{The arity-one equation: solutions and free directions}

\begin{theorem}[The arity-one equation]
\label{thm:c12-eq1}
Let $\hbar F_1^{(0)}$ solve the arity-zero equation of \eqref{eq:c12-eqs} and let $x$ be a generator.
\begin{itemize}
\item (i) $\mathcal R^{(1)}(x)$ is $(Q+\{I,-\})$-closed, of degree $0$, Weyl weight $4-\tfrac12\lambda_x$ and position degree at most $p_x$, with the covariance of Theorem \ref{thm:c12-bound}.
\item (ii) The equation $(Q+\{I,-\})\hbar F_1^{(1)}(x)=\mathcal R^{(1)}(x)$ has a local solution exactly when the class $[\mathcal R^{(1)}(x)]$ vanishes in the cohomology $H^0(Q+\{I,-\})$ of the local functionals with these quantum numbers.
\item (iii) Two solutions differ by a $(Q+\{I,-\})$-closed local functional of degree $-1$ with the same quantum numbers, and modulo exact ones by a class of $H^{-1}(Q+\{I,-\})$.
\item (iv) Replacing $\hbar F_1^{(0)}$ by $\hbar F_1^{(0)}+G^{(0)}$ with $(Q+\{I,-\})G^{(0)}=0$ replaces $\mathcal R^{(1)}(x)$ by $\mathcal R^{(1)}(x)-\{G^{(0)},S^{(1)}(x)\}$; for $G^{(0)}=(Q+\{I,-\})\phi^{(0)}$ the change is absorbed by $\hbar F_1^{(1)}(x)\mapsto\hbar F_1^{(1)}(x)+\{\phi^{(0)},S^{(1)}(x)\}$, the field redefinition generated by $\phi^{(0)}$ (Proposition \ref{prop:c8-laws} (a)).
\item (v) In $H^0(Q+\{I,-\})$,
\begin{equation}
[\mathcal R^{(1)}([x,y])]=(-1)^{|x|}[\{\mathcal R^{(1)}(x),S^{(1)}(y)\}]-[\{S^{(1)}(x),\mathcal R^{(1)}(y)\}],
\label{eq:c12-cocycle1}
\end{equation}
the Chevalley--Eilenberg one-cocycle condition of $\mathfrak{psl}(4|4)$ with values in $H^0(Q+\{I,-\})$, the action being the bracket with $S^{(1)}(y)$; the change (iv) is the coboundary of $[G^{(0)}]$, and the obstruction to solving through arity one, for some arity-zero solution, is the class of $x\mapsto[\mathcal R^{(1)}(x)]$ in $H^1(\mathfrak{psl}(4|4);H^0(Q+\{I,-\}))$.
\end{itemize}
\end{theorem}

\begin{proof}
(i) Lemma \ref{lem:c12-wz} (iii). (ii), (iii) Definition of the cohomology of local functionals. (iv) $\{(Q+\{I,-\})\phi^{(0)},S^{(1)}(x)\}=-(Q+\{I,-\})\{\phi^{(0)},S^{(1)}(x)\}$ by \eqref{eq:c12-deriv}. (v) The arity-two identity of Lemma \ref{lem:c12-wz} (ii) for $G=\hbar\,\mathrm{Ob}-s(\hbar F_1)$, whose arity-zero component vanishes, reads $(Q+\{I,-\})G^{(2)}(x,y)=G^{(1)}([x,y])-(-1)^{|x|}\{G^{(1)}(x),S^{(1)}(y)\}+\{S^{(1)}(x),G^{(1)}(y)\}$ with $G^{(1)}=\mathcal R^{(1)}$; in cohomology the left side vanishes. For $\mathcal R^{(1)}(x)=-\{G^{(0)},S^{(1)}(x)\}$ the identity follows from $(-1)^{|x|}\{S^{(1)}(x),S^{(1)}(y)\}=-S^{(1)}([x,y])-(Q+\{I,-\})S^{(2)}(x,y)$ (Lemma \ref{lem:c2-components}) and \eqref{eq:c12-deriv} for $\{G^{(0)},-\}$, the term $\{G^{(0)},(Q+\{I,-\})S^{(2)}(x,y)\}$ being exact.
\end{proof}

On two legs the bracket with $I$ raises the number of legs, so the equation reads $Q\,\hbar F_1^{(1)}(x)|_2=\mathcal R^{(1)}(x)|_2$, and the relevant cohomology is that of $Q$ on quadratic functionals.

\begin{proposition}[Two legs along $\Sigma$]
\label{prop:c12-xi2}
Let $\xi\in S_+\otimes W$; for $\xi\in S_-\otimes W^*$ exchange the chiralities. With $\Gamma_V:=\Gamma(\xi,\lambda)|_V$, $\Gamma_\perp:=\Gamma(\xi,\lambda)|_{W_{\mathbb R}^\perp}$ and $(\xi,\lambda^*)$ the $\mathfrak g$-valued pairings, define
\[
g_1:=\int(\Gamma_V,d^*d\,A|_V),\quad g_2:=\int(\Gamma_V,dd^*A|_V),\quad g_3:=\int(\Gamma_\perp,d^*d\,A|_{W_{\mathbb R}^\perp}),\quad g_4:=\int\big(d^*dc,(\xi,\lambda^*)\big),\quad T_0:=\int\big(d^*A|_V,(\xi,\lambda^*)\big).
\]
\begin{itemize}
\item (i) The quadratic covariant local functionals of degree $-1$ and Weyl weight $\tfrac92$, linear in $\xi$, are spanned by $T_0,T_1,T_2,T_3,T_4$; those of degree $0$ by $g_1,g_2,g_3,g_4$.
\item (ii) $QT_1=-QT_3=-\tfrac12g_1$, $QT_2=-QT_4=-\tfrac12g_3$ and $QT_0=-(g_2+g_4)$. The kernel of $Q$ at degree $-1$ is spanned by $S^V_2(\xi)=T_1+T_3$ and $S^\perp_2(\xi)=T_2+T_4$, and no nonzero combination of them is $Q$-exact.
\item (iii) The $Q$-closed functionals of degree $0$ are spanned by $g_1$, $g_3$ and $g_2+g_4$, which is the image of $Q$: the cohomology $H^0(Q)$ of these quadratic functionals is zero.
\item (iv) The two-leg part of the arity-one equation along $\xi$ is solvable for every arity-zero solution: with $\mathcal R^{(1)}(\xi)|_2=r_1g_1+r_3g_3+r_0(g_2+g_4)$,
\begin{equation}
\hbar F_1^{(1)}(\xi)|_2=-2r_1T_1-2r_3T_2-r_0T_0+\theta^\xi_V\,S^V_2(\xi)\ \text{units}+\theta^\xi_\perp\,S^\perp_2(\xi)\ \text{units},
\label{eq:c12-xisol}
\end{equation}
with free $\theta^\xi_V$, $\theta^\xi_\perp$. The coefficient $r_0$ of the gauge-fixing structure receives the finite counterterm $\tfrac{27}{16}\int(d^*A|_V,d^*A|_V)$ units through $\{\int(d^*A|_V,d^*A|_V),T_1\}=2g_2$, and the parts $g_2$, $g_4$ of $\hbar\,\mathrm{Ob}^{(1)}(\xi)|_2$.
\item (v) On two legs the Wess--Zumino identity is $Q\,\mathrm{Ob}^{(1)}(\xi)|_2=-\{\mathrm{Ob}^{(0)}|_2,T_1\}$, with $\hbar\,\mathrm{Ob}^{(0)}|_2=-\tfrac{27}8\int(d^*dc,d^*A|_V)$ units (Theorem \ref{thm:c4-obstruction}): it fixes the components of $\mathrm{Ob}^{(1)}(\xi)|_2$ along $g_2$, $g_4$ outside the closed combination.
\item (vi) The field-counting directions $(Q+\{I,-\})\mathcal N_f$ of $\hbar F_1^{(0)}$ (Theorem \ref{thm:c9-F3}) shift $\hbar F_1^{(1)}(\xi)|_2$ by $\{\mathcal N_f,S^{(1)}_2(\xi)\}$, which multiplies each arrow $f\to f'$ of the classical transformation by the difference of the counting weights (Theorem \ref{thm:c12-eq1} (iv)).
\end{itemize}
\end{proposition}

\begin{proof}
(i) $\xi\otimes\lambda|_{S_-\otimes W^*}\in(S_+\otimes S_-)\otimes(W\otimes W^*)$ contains the $SU(4)$ singlet in $V\otimes\mathbf 1$, so every bilinear in $(\xi,\lambda|_{S_-\otimes W^*})$ is a component of $\Gamma_V$; with $A|_V$ and two derivatives, after integration by parts onto $A|_V$, the contractions are $\Gamma_V\cdot d^*dA|_V$ and $\Gamma_V\cdot dd^*A|_V$. $\xi\otimes\lambda|_{S_+\otimes W}\in(\mathbf 1\oplus\wedge^2_+)\otimes(\mathbf 6\oplus\mathbf{10})$ pairs with $A|_{W_{\mathbb R}^\perp}$ through $\Gamma_\perp$, and $\xi\otimes\lambda^*|_{S_+\otimes W^*}$ contains the scalar $(\xi,\lambda^*)$, paired with $c$ and two derivatives. At degree $-1$, $(\lambda,A^*)$ without derivatives gives $T_1$, $T_2$, and $(A,\lambda^*)$ with one derivative gives $T_0$, $T_3$, $T_4$ (Proposition \ref{prop:c11-freexi} (i)). (ii) $Q$ substitutes $A^*\mapsto-\tfrac12d^*dA$, $\lambda^*\mapsto{\not d}\lambda$, $A|_V\mapsto-dc$ (Proposition \ref{prop:c1-Q}): $T_1\mapsto-\tfrac12\Gamma_V\cdot d^*dA|_V$, $T_2\mapsto-\tfrac12\Gamma_\perp\cdot d^*dA|_{W_{\mathbb R}^\perp}$; in $T_0$ the substitution in $A|_V$ gives $-g_4$ and that in $\lambda^*$ gives $(\xi,{\not d}\lambda)\,d^*A|_V$ with $(\xi,{\not d}\lambda)=-d^*\Gamma_V$, hence $-g_2$. $S^V_2$ and $S^\perp_2$ are closed, the two sector parts of the quadratic part of $\{S,S^{(1)}(\xi)\}=0$. The images $g_1$, $g_3$, $g_2+g_4$ are independent, so the kernel is two-dimensional; a $Q$-exact functional of degree $-1$ would be $Q$ of a quadratic functional of degree $-2$ and weight $\tfrac92$, and these are zero (Proposition \ref{prop:c11-freexi} (i)). (iii) $Qg_1=Qg_3=0$, and $Qg_2=-\int(\Gamma_V,dd^*dc)=-Qg_4$, nonzero; the closed space is three-dimensional and equals the image. (iv) The residual is closed (Lemma \ref{lem:c12-wz} (iii)); (ii) and (iii) give the solution, and $\{\int(d^*A|_V,d^*A|_V),T_1\}$ contracts $A|_V$ with $A^*|_V$. (v) Lemma \ref{lem:c12-wz} (iii) on two legs. (vi) The quadratic part of $(Q+\{I,-\})\mathcal N_f$ is $Q\mathcal N_f$, and $\{Q\mathcal N_f,S^{(1)}_2(\xi)\}=-Q\{\mathcal N_f,S^{(1)}_2(\xi)\}$ by \eqref{eq:c12-deriv}.
\end{proof}

\begin{proposition}[Two legs along the even generators, $\Sigma^*$ and $K_b$]
\label{prop:c12-two}
\begin{itemize}
\item (i) $\Delta$. The $Q$-closed quadratic functionals of degree $-1$, Weyl weight $4$, position degree at most one and reflection-even are spanned, modulo $Q$-exact ones, by the three sector parts of $S^{(1)}(\Delta)$; the exact ones by $Q\int(\iota_xA|_V,c^*)$ and $Q\int(\rho(x)\lambda^*,\lambda^*)$ (Proposition \ref{prop:c11-tc0}). $H^0(Q)$ vanishes on the two-leg target.
\item (ii) $P_b$. The translation-invariant quadratic closed functionals of degree $-1$, weight $5$, linear in $b$, are spanned modulo exact ones by the three sector parts of $S^{(1)}(P_b)$, the position parts of those of (i); the exact ones by $Q\int(\iota_bA|_V,c^*)$ and $Q\int(\rho(b)\lambda^*,\lambda^*)$.
\item (iii) $M_\omega$, $R_r$. The sector parts are closed. Along $M_\omega$, in the leg sets $(A|_V,A|_V)_2$, $(c,A^*|_V)_1$ the closed functionals of degree $0$ are spanned by $Q\int\omega_{\mu\nu}(A_\mu,A^*_\nu)$, and in the leg set $(\lambda,\lambda)_1$ by $Q\int(\tfrac14\omega_{\mu\nu}\rho_\mu\rho_\nu\lambda,\lambda^*)$, the structure $\int(\lambda,\omega_{\mu\nu}\rho_{[\mu}\rho_\nu\rho_{\kappa]}\partial_\kappa\lambda)$ being a total derivative: $H^0(Q)=0$ there, and the spin term
\[
S^{(1)}(M_\omega)_{\rm spin}:=\int\omega_{\mu\nu}(A_\mu,A^*_\nu)+\int\big(\tfrac14\omega_{\mu\nu}\rho_\mu\rho_\nu\lambda,\lambda^*\big)
\]
of $\hbar F_1^{(1)}(M_\omega)$ is fixed by $\mathrm{Ob}^{(1)}(M_\omega)|_2$. Along $R_r$, $\hbar F_1^{(1)}(R_r)=0$ solves the equation (Theorem \ref{thm:c12-second} (iii)).
\item (iv) $\eta\in S_+\otimes W^*$. The translation-invariant quadratic functionals of degree $-1$ and weight $\tfrac72$ are spanned by $\int(\rho(A|_V)\eta,\lambda^*)$ and $\int s_0$; $Q$ is block triangular for the split into position parts and intrinsic parts; its kernel on the full degree-$(-1)$ space is spanned by $S^V(\eta)$ and $S^\perp(\eta)$; the cohomology $H^0(Q)$ on the two-leg degree-$0$ functionals of weight $\tfrac72$ with position degree at most one is one-dimensional, a class in the sector $(A|_{W_{\mathbb R}^\perp},\lambda)$.
\item (v) $K_b$. The sector parts $S^{(1)}(K_b)_d$ are closed; with position degree at most two, $H^{-1}(Q)$ and $H^0(Q)$ on two legs have dimensions $7$ and $4$: $H^{-1}$ is spanned by the three classes $[S^{(1)}(K_b)_d]$ and four second-order symmetries of the free fields, written $G^K_1,\dots,G^K_4$, closed and outside the span of the $S^{(1)}(z)$. $\{\mathrm{Ob}^{(0)}|_2,S^{(1)}(K_b)\}$ and $\{\hbar F_1^{(0)},S^{(1)}(K_b)\}$ are nonzero, $\int(d^*dc,d^*A|_V)$ being outside the special conformal invariants.
\end{itemize}
\end{proposition}

\begin{proof}
(i) Proposition \ref{prop:c11-tc0} and Theorem \ref{thm:c9-quadratic} (iii). (ii) The position part along $b$ commutes with $Q$ (Lemma \ref{lem:c11-canonical}) and maps the closed and exact directions of (i) to those listed. (iii) $v.$ commutes with $Q$; in the vector sector $Q\int\omega_{\mu\nu}(\partial_\mu A_\nu,d^*A|_V)$ and $Q\int\omega_{\mu\nu}(\partial_\mu c,A^*_\nu)$ are proportional and nonzero, so the closed subspace of these two is one-dimensional and contains $Q\int\omega_{\mu\nu}(A_\mu,A^*_\nu)\neq0$; in the gaugino sector $\Gamma$ is symmetric. (iv), (v) The leg sets of Table \ref{tab:c12-legs} and their position-dependent completions, with the exact ranks of $Q$ on plane-wave normal forms; $Q$ commutes with $-\{S^{(1)}(P_b),-\}$, which gives the triangular form. $\{\int(d^*dc,d^*A|_V),S^{(1)}(K_b)\}$ is the variation of a functional outside the special conformal invariants.
\end{proof}

\begin{theorem}[The arity-one obstructions on two legs]
\label{thm:c12-ob1}
On two legs the arity-one obstructions are
\begin{equation}
\begin{aligned}
\hbar\,\mathrm{Ob}^{(0)}|_2&=-\tfrac{27}8\int(d^*dc,d^*A|_V)\ \text{units},\\
\hbar\,\mathrm{Ob}^{(1)}(\Delta)|_2&=\Big[-\tfrac{11}3\int(dA,dA)+\tfrac{47}3\int(\lambda,{\not d}\lambda)-\tfrac{14}3\int(dc,A^*)\Big]\ \text{units},\\
\hbar\,\mathrm{Ob}^{(1)}(P_b)|_2&=0,\qquad \hbar\,\mathrm{Ob}^{(1)}(R_r)|_2=0,\qquad \hbar\,\mathrm{Ob}^{(1)}(M_\omega)|_2=\tfrac13\,Q\,S^{(1)}(M_\omega)_{\rm spin}\ \text{units},\\
\hbar\,\mathrm{Ob}^{(1)}(\xi)|_2&=\Big[\big(\tfrac{79}{24}-\tfrac12\log\tfrac32\big)(g_1+g_3)+\tfrac{27}8\,g_2+0\cdot g_4\Big]\ \text{units},\\
\hbar\,\mathrm{Ob}^{(1)}(\eta)|_2&=\sum_\mu x^\mu\,\hbar\,\mathrm{Ob}^{(1)}(\rho(e_\mu)\eta)|_2+\Big[-\tfrac52\int(\eta,\lambda)\,d^*A|_V-\tfrac52\int(\eta,\rho(dA|_V)\lambda)\\
&\qquad+\big(-\tfrac32+3\log\tfrac32\big)\int(\eta,\rho(dA|_{W_{\mathbb R}^\perp})\lambda)+\tfrac52\int(\rho(dc)\eta,\lambda^*)\Big]\ \text{units}+Q(\cdot),
\end{aligned}
\label{eq:c12-ob1}
\end{equation}
the same for both chiralities of $\xi$ and $\eta$, the intrinsic part along $\eta$ in the canonical density of Definition \ref{def:c11-canonical}. Along $\Delta$ the coefficient of $\int(d^*A|_V,d^*A|_V)$ and of every functional with an explicit factor $x$ is $0$; along $K_b$ the obstruction lies in the span of the canonical data with position degree at most two. In every case the poles of the defect cycles sum to $\tfrac94\{\int(A,A),S^{(1)}(x)\}$ units, as \eqref{eq:c12-pole} requires; along $\Delta$ this is $-9\int(A,A)$ units and along $\xi$ it is $\tfrac92\big[\int(\Gamma_V,A|_V)+\int(\Gamma_\perp,A|_{W_{\mathbb R}^\perp})\big]$ units. The logarithmic coefficient along $\xi$ is $L^{(1)}(\xi)|_2=\{\mathcal N,S^{(1)}_2(\xi)\}$ units $=2(T_1+T_2-T_3-T_4)$ units, so the closed combination \eqref{eq:c12-closedlog} vanishes on two legs along $\xi$. The Wess--Zumino identities \eqref{eq:c12-wz} hold on every entry.
\end{theorem}

\begin{proof}
The cycles of Theorem \ref{thm:c12-classes}, with the placements of Table \ref{tab:c12-place}, evaluated by \eqref{eq:c12-special} and the closed forms of Proposition \ref{prop:c12-closed}, decomposed on the bases of Propositions \ref{prop:c12-xi2} and \ref{prop:c12-two}; the position part along $\eta$ is Proposition \ref{prop:c12-pos}. Along $\Delta$ the value is the target of Theorem \ref{thm:c9-quadratic} (ii), here recomputed through the defect cycles; along $P_b$ its vanishing is the position part of the statement along $\Delta$ (Proposition \ref{prop:c12-pos}). $L^{(1)}(\xi)|_2$ is the coefficient of $\log\varepsilon$ of the composite and cubic cycles, and $\{\mathcal N,-\}$ multiplies the arrows of $T_1$, $T_2$ ($\lambda\to A^*$) by $-2+4=2$ and those of $T_3$, $T_4$ ($A\to\lambda^*$) by $2-4=-2$.
\end{proof}

\begin{theorem}[Two-leg solutions with all free directions]
\label{thm:c12-sol}
On two legs the arity-zero solution is $\hbar F_1^{(0)}|_2=\tfrac{27}{16}\int(d^*A|_V,d^*A|_V)$ units plus the $s$-closed arity-zero directions, and the arity-one solutions are
\begin{equation}
\begin{aligned}
\hbar F_1^{(1)}(\Delta)|_2&=\Big[\tfrac{22}3\int(A,A^*)-\tfrac{47}3\int(\lambda^*,\lambda)-\tfrac83\int(c^*,c)\Big]\ \text{units}+\sum_d\theta^\Delta_d\,S^{(1)}(\Delta)_d\ \text{units}\\
&\quad+\kappa_V\,Q\int(\iota_xA|_V,c^*)\ \text{units}+\kappa_\lambda\,Q\int(\rho(x)\lambda^*,\lambda^*)\ \text{units},\\
\hbar F_1^{(1)}(P_b)|_2&=\sum_d\theta^{P_b}_d\,S^{(1)}(P_b)_d\ \text{units}+(\text{the parts along }b\text{ of the }\kappa\text{ directions}),\\
\hbar F_1^{(1)}(M_\omega)|_2&=\tfrac13\,S^{(1)}(M_\omega)_{\rm spin}\ \text{units}+\sum_d\theta^{M}_d\,S^{(1)}(M_\omega)_d\ \text{units},\qquad \hbar F_1^{(1)}(R_r)|_2=\sum_d\theta^R_d\,S^{(1)}(R_r)_d\ \text{units},\\
\hbar F_1^{(1)}(\xi)|_2&=\big(\log\tfrac32-\tfrac{79}{12}\big)(T_1+T_2)\ \text{units}+\theta^\xi_V\,S^V_2(\xi)\ \text{units}+\theta^\xi_\perp\,S^\perp_2(\xi)\ \text{units},\\
\hbar F_1^{(1)}(\eta)|_2&=\sum_\mu x^\mu\big(\log\tfrac32-\tfrac{79}{12}\big)(T_1+T_2)(\rho(e_\mu)\eta)\ \text{units}-\big(3\log\tfrac32-\tfrac32\big)\int s_0\ \text{units}-\tfrac52\int(\rho(A|_V)\eta,\lambda^*)\ \text{units}\\
&\quad+\theta^\eta_V\,S^V(\eta)\ \text{units}+\theta^\eta_\perp\,S^\perp(\eta)\ \text{units},\\
\hbar F_1^{(1)}(K_b)|_2&=\int(b,x)\Big\{-\tfrac{22}3(A,A^*)+16(\lambda,\lambda^*)+\tfrac83(c,c^*)\Big\}\ \text{units}+\tfrac23\int\big\{(\iota_bA)(\iota_xA^*)-(\iota_xA)(\iota_bA^*)\big\}\ \text{units}\\
&\quad-\tfrac13\int b_\mu x_\nu(\lambda,\rho_\mu\rho_\nu\lambda^*)\ \text{units}+\sum_d\theta^K_d\,S^{(1)}(K_b)_d\ \text{units}+\sum_{j=1}^4\kappa^K_jG^K_j,
\end{aligned}
\label{eq:c12-sol}
\end{equation}
with the free directions displayed: the sector coefficients $\theta^x_d$, the exact directions $\kappa_V$, $\kappa_\lambda$, the four closed directions $\kappa^K_j$ along $K_b$, and, along $\xi$, the shifts of Proposition \ref{prop:c12-xi2} (vi). $Q$ of the first bracket of $\hbar F_1^{(1)}(\Delta)|_2$ is $\hbar\,\mathrm{Ob}^{(1)}(\Delta)|_2$. The part of $\hbar F_1^{(1)}(K_b)|_2$ linear in $x$ is $-(b,x)$ times the particular part along $\Delta$ plus the spin term of the rotation part of $\partial_{b'}v$, $16-\tfrac13=\tfrac{47}3$ on the gaugino. On two legs the residual has no component along $H^0(Q)$ in each case: $H^0(Q)$ vanishes along $\Delta$, $P_b$, $M_\omega$, $R_r$, $\Sigma$, and along $\Sigma^*$ and $K_b$ the computed residual is orthogonal to its classes.
\end{theorem}

\begin{proof}
Propositions \ref{prop:c12-xi2}, \ref{prop:c12-two} applied to the values of Theorem \ref{thm:c12-ob1}, with the residuals of Definition \ref{def:c12-ob}. Along $\xi$, $\hbar\,\mathrm{Ob}^{(1)}(\xi)|_2-\{\hbar F_1^{(0)}|_2,T_1\}$ has $r_0=0$: the coefficient $\tfrac{27}8$ of $g_2$ in the obstruction equals the contribution $\tfrac{27}{16}\cdot2$ of the finite counterterm, and $g_4$ has coefficient $0$; so $\hbar F_1^{(1)}(\xi)|_2$ has no term $T_0$, and $-2r_1=-2r_3=\log\tfrac32-\tfrac{79}{12}$. Along $\eta$ the particular part is read on the full degree-$(-1)$ basis; along $K_b$ the $x$-linear part follows from Proposition \ref{prop:c12-pos} with the rotation part of $\partial_{b'}v$.
\end{proof}

\begin{theorem}[The cubic part fixes the sector differences]
\label{thm:c12-cubic}
Write the quadratic part of a solution as $\hbar F_1^{(1)}(x)|_2=F^{\rm part}_2(x)+\sum_d\theta^x_dS^{(1)}(x)_d$ units $+$ exact terms, with $I_3$ the cubic part of $I$. The part of the arity-one equation with three legs is
\[
Q\,\hbar F_1^{(1)}(x)|_3+\{I_3,\hbar F_1^{(1)}(x)|_2\}=\mathcal R^{(1)}(x)|_3,\qquad \mathcal R^{(1)}(x)|_3=\hbar\,\mathrm{Ob}^{(1)}(x)|_3-\{\hbar F_1^{(0)}|_2,S^{(1)}_3(x)\}-\{\hbar F_1^{(0)}|_3,S^{(1)}_2(x)\}.
\]
\begin{itemize}
\item (i) Each $\{I_3,S^{(1)}(x)_d\}$ is $Q$-closed and $\sum_d\{I_3,S^{(1)}(x)_d\}=-QS^{(1)}_3(x)$; the equation depends on the $\theta^x_d$ through $\sum_{d\neq d_0}(\theta^x_d-\theta^x_{d_0})[\{I_3,S^{(1)}(x)_d\}]$ in the cohomology $H^0(Q)_3$ of the cubic functionals of degree $0$, weight $4-\tfrac12\lambda_x$ and position degree at most $p_x$.
\item (ii) It is solvable exactly when $[\mathcal R^{(1)}(x)|_3-\{I_3,F^{\rm part}_2(x)\}]=\sum_{d\neq d_0}(\theta^x_d-\theta^x_{d_0})[\{I_3,S^{(1)}(x)_d\}]$ units in $H^0(Q)_3$.
\item (iii) The classes $[\{I_3,S^{(1)}(x)_d\}]$, $d\neq d_0$, are linearly independent exactly when no combination $\sum_dc_dS^{(1)}(x)_d$ other than the multiples of $S^{(1)}_2(x)$ extends to a $(Q+\{I,-\})$-closed functional through three legs; then the cubic equation fixes the differences $\theta^x_d-\theta^x_{d_0}$.
\end{itemize}
\end{theorem}

\begin{proof}
On three legs $(Q+\{I,-\})$ applied to a functional with two and three legs is $Q$ of the cubic part plus the bracket of $I_3$ with the quadratic part. (i) $QI_3=0$ and $QS^{(1)}(x)_d=0$ give $Q\{I_3,S^{(1)}(x)_d\}=0$ by \eqref{eq:c12-deriv}; the cubic part of $\{S,S^{(1)}(x)\}=0$ is $QS^{(1)}_3(x)+\{I_3,S^{(1)}_2(x)\}=0$. (ii) Definition of $H^0(Q)_3$; the exact quadratic directions contribute $\{I_3,Q\psi\}=\pm Q\{I_3,\psi\}$. (iii) A relation $\sum_dc_d[\{I_3,S^{(1)}(x)_d\}]=0$ means $\{I_3,\sum_dc_dS^{(1)}(x)_d\}=-QY_3$, that is, $(Q+\{I,-\})(\sum_dc_dS^{(1)}(x)_d+Y_3)=0$ through three legs; $\sum_dS^{(1)}(x)_d$ always extends, by $Y_3=S^{(1)}_3(x)$.
\end{proof}

\begin{proposition}[The arity-one equations along $\Sigma$ constrain the arity-zero couplings]
\label{prop:c12-couplings}
Let $\mathcal H_4$ be $H^0(Q+\{I,-\})$ on the colour-singlet local functionals of degree $0$ and weight $4$, invariant under $\mathrm{Spin}(4)\times\mathrm{Spin}(6)$ and the reflection: the classical couplings modulo field redefinitions, containing the class of $S$ and the classes of Theorem \ref{thm:c9-F3} that extend to closed functionals. For $\xi\in\Sigma$ let $\delta_\xi[G]:=[\{G,S^{(1)}(\xi)\}]$. The arity-one equations along $\Sigma$ are solvable for some arity-zero solution exactly when $[\mathcal R^{(1)}(\xi)]$ lies in the image of $\delta_\xi$ for every $\xi$, one $[G]$ serving all $\xi$; the arity-zero freedom left after them is $\bigcap_\xi\ker\delta_\xi$, which contains $[S]$. Along $\Delta$, $P_b$, $M_\omega$, $R_r$ and $K_b$ the map $[G]\mapsto[\{G,S^{(1)}(x)\}]$ vanishes on $\mathcal H_4$.
\end{proposition}

\begin{proof}
Theorem \ref{thm:c12-eq1} (iv), (v). For even $x$, $\{G,S^{(1)}(x)\}=\mp v.G$ (Lemma \ref{lem:c9-SDelta} (i), Definition \ref{def:c7-vdot}), zero for a translation- and rotation-invariant $G$ of weight $4$; along $K_b$ the gauge-invariant functionals of weight $4$ of $A$ and $\lambda$ are conformally invariant with the weights $1$, $\tfrac32$ (Proposition \ref{prop:c2-coeff}), and the antifield terms of weight $4$ because $S$ is.
\end{proof}

\begin{definition}[The classical symmetry condition]
\label{def:c12-cs}
The classical symmetry condition at the weight of $x$ (resp. of a pair $(x,y)$) is: every $(Q+\{I,-\})$-closed colour-singlet local functional of degree $-1$, of Weyl weight $4-\tfrac12\lambda_x$ (resp. $4-\tfrac12(\lambda_x+\lambda_y)$), position degree at most $p_x$ (resp. $p_x+p_y$), linear in the Killing data and equivariant as in Theorem \ref{thm:c12-bound}, is a linear combination of the functionals $S^{(1)}(z)$, $z\in\mathfrak{psl}(4|4)$, with these quantum numbers, plus a $(Q+\{I,-\})$-exact functional. It states that the global symmetries of the interacting classical theory at these weights are the superconformal ones. On two legs $H^{-1}(Q)$ is larger, containing the sector parts $S^{(1)}(z)_d$ and, at the weights with two powers of the position, second-order symmetries of the free fields built from conformal Killing tensors; the condition concerns their extensions through three legs, which Theorem \ref{thm:c12-cubic} (iii) tests.
\end{definition}

\begin{theorem}[Solution space at arity one]
\label{thm:c12-H}
Assume the classical symmetry condition at the weights of the generators and that the arity-one equations are solvable. Then for every generator $x$ the solutions are $\hbar F_1^{(1)}(x)=F^{\rm part}(x)+t_x\,S^{(1)}(x)$ units $+(Q+\{I,-\})(\text{local})$, with one number $t_x$, the sector coefficients being equal to $t_x$ up to the fixed differences of Theorem \ref{thm:c12-cubic} contained in $F^{\rm part}(x)$. In a scheme equivariant for translations, rotations and $\mathfrak{so}(W_{\mathbb R}^\perp)$, $t_x$ depends on the block of $x$ alone: $t_\Delta$, $t_P$, $t_M$, $t_R$, $t_{\xi_+}$, $t_{\xi_-}$, $t_{\eta_+}$, $t_{\eta_-}$, $t_K$.
\end{theorem}

\begin{proof}
Theorem \ref{thm:c12-eq1} (iii) and the condition of Definition \ref{def:c12-cs}: the functional $S^{(1)}(z)$ with the weight, position degree and Killing data of $x$ that is equivariant is $S^{(1)}(x)$ itself. Equivariance under the compact groups makes $t_x$ constant on each irreducible block.
\end{proof}

\begin{remark}[The free directions and the equations that fix them]
\label{rem:c12-freetable}
The differences of the sector coefficients are fixed by the cubic part of the arity-one equation (Theorem \ref{thm:c12-cubic}) and, along $\Sigma$, $\Sigma^*$, $P_b$, $\Delta$, $M_\omega$, $R_r$, already by the two-leg parts of the arity-two equations (Theorem \ref{thm:c12-ar2}); the uniform parts by the arity-two equations of the next section:

\begin{tabular}{l|l|l}
$x$ & closed quadratic directions of $\hbar F_1^{(1)}(x)$ & uniform part fixed by\\
\hline
$\Delta$ & $S^{(1)}(\Delta)_d$; $Q\int(\iota_xA|_V,c^*)$, $Q\int(\rho(x)\lambda^*,\lambda^*)$ & $(\Delta,P_b)$, $(\Delta,\xi)$, $(\xi,\eta)$\\
$P_b$ & $S^{(1)}(P_b)_d$; $Q\int(\iota_bA|_V,c^*)$, $Q\int(\rho(b)\lambda^*,\lambda^*)$ & up to the inner direction\\
$M_\omega$ & $S^{(1)}(M_\omega)_d$ & $(K_b,P_{b'})$, $(\xi,\eta)$\\
$R_r$ & $S^{(1)}(R_r)_\perp$, $S^{(1)}(R_r)_\lambda$ & $(\xi,\eta)$\\
$\Sigma$ & $S^V_2(\xi)$, $S^\perp_2(\xi)$; difference fixed by $(\xi_\pm,\xi'_\pm)$ & up to the inner and outer directions\\
$\Sigma^*$ & $S^V(\eta)$, $S^\perp(\eta)$; difference fixed by $(\xi_\pm,\eta_\mp)$ & up to the inner and outer directions\\
$K_b$ & $S^{(1)}(K_b)_d$, $G^K_1,\dots,G^K_4$ & up to the inner direction\\
\end{tabular}
\tabcaption{The free directions of the arity-one solutions and the equations that fix them; the inner and outer directions are those of Theorem \ref{thm:c12-fix} (ii).\label{tab:c12-free}}
\end{remark}

\section{Arity two: the coefficients and the one-loop algebra}

\begin{proposition}[Graph classes at arity two]
\label{prop:c12-classes2}
Fix generators $x$, $y$ and $n\ge2$. At order $\hbar$, with $F_1=0$, the part with $n$ legs of the arity-two component along $(x,y)$ of the flow of $\mathfrak d$ consists of: (C2) irreducible cycles with exactly one special edge and either two vertices $S^{(1)}(x)$, $S^{(1)}(y)$ or one vertex $S^{(2)}(x,y)$ on the cycle, the other vertices from $I$; (Rd2) reducible graphs whose subgraphs are cycles of type (C) along $x$ or $y$, of type (C2) with fewer legs, or of the arity-zero defect, joined by edges $P(\varepsilon,T)$ to trees; (Br2) bridges; (Ct2) the counterterm trees of $\mathfrak d^{(2)}(x,y)$, of $\mathfrak d^{(1)}$ with a separate $S^{(1)}(y)$ or $S^{(1)}(x)$, and of $\mathfrak d^{(0)}$ with two separate vertices. Two vertices $S^{(1)}(x)$, $S^{(1)}(y)$ in the line placement are adjacent through an $A|_V\otimes A|_V$ edge when both are even (the chain $c$--$A|_V$--$A|_V$--$c$, a double composite line) or through a $\lambda\otimes\lambda$ edge when both are odd, and never when one is even and one odd. As $\varepsilon\to0$, (Br2) with the special-edge parts of (Ct2) tends to $0$, (Rd2) tends to the trees of the components of $\mathrm{Ob}$ of lower arity or with fewer legs, the $\log\varepsilon$ terms of (Ct2) cancel, (Ct2) has no pole, $\{\int(A,A),S^{(2)}(x,y)\}$ being $0$, and
\begin{equation}
\sum_{\mathcal C\in(C2),\,n}a_m(\mathcal C)=0\ (m\ge1),\qquad \hbar\,\mathrm{Ob}^{(2)}(x,y)\big|_{n\ \text{legs}}=\sum_{\mathcal C\in(C2),\,n\ \text{legs}}a_0(\mathcal C).
\label{eq:c12-local2}
\end{equation}
For two odd generators the arity-two generating cochain is $I+c\,S^{(1)}(\xi)+c'S^{(1)}(\xi')-cc'S^{(2)}(\xi,\xi')$ with odd parameters $c$, $c'$, the sign forced by the classical master equation in the form $QS^{(2)}_2(\xi,\xi')+S^{(1)}_2(P_{\Gamma(\xi,\xi')|_V})-\{S^{(1)}_2(\xi),S^{(1)}_2(\xi')\}=0$.
\end{proposition}

\begin{proof}
As Theorem \ref{thm:c12-classes}, with Lemma \ref{lem:c12-defect} for $\mathfrak d^{(2)}$: its marked vertices of order $\hbar^0$ are the special-edge parts of $(-1)^{|x|}\{S^{(1)}(x),S^{(1)}(y)\}_\varepsilon$ and $\{I,S^{(2)}(x,y)\}_\varepsilon$, and the marked vertices of lower arity enter with the remaining generators on separate vertices. $\int(A,A)$ has no antifield and $S^{(2)}$ no leg $A^*$. The adjacency rules follow from the components of $P$ (Proposition \ref{prop:c3-P}) and the legs of Table \ref{tab:c12-place}. The generating cochain: the arity-two component of $d_{CE}f+\tfrac12\{f,f\}$ at the two odd parameters, with Lemma \ref{lem:c2-components}.
\end{proof}

\begin{theorem}[The arity-two equation and its cocycle]
\label{thm:c12-eq2}
Let $\hbar F_1^{(0)}$, $\hbar F_1^{(1)}$ solve the equations of arity zero and one.
\begin{itemize}
\item (i) $\mathcal R^{(2)}(x,y)$ is $(Q+\{I,-\})$-closed of degree $-1$, Weyl weight $4-\tfrac12(\lambda_x+\lambda_y)$ and position degree at most $p_x+p_y$; the arity-two equation along $(x,y)$ is solvable exactly when $[\mathcal R^{(2)}(x,y)]=0$ in $H^{-1}(Q+\{I,-\})$.
\item (ii) Replacing $\hbar F_1^{(1)}$ by $\hbar F_1^{(1)}+\varphi$ with $(Q+\{I,-\})\varphi(x)=0$ for all $x$ replaces $\mathcal R^{(2)}(x,y)$ by $\mathcal R^{(2)}(x,y)-\varphi([x,y])-(-1)^{|x|}\big(\{\varphi(x),S^{(1)}(y)\}+\{S^{(1)}(x),\varphi(y)\}\big)$, minus the Chevalley--Eilenberg coboundary of $[\varphi]$ with values in $H^{-1}(Q+\{I,-\})$; exact $\varphi$ change $\mathcal R^{(2)}$ by exact terms.
\item (iii) $\{\hbar F_1^{(1)}(x),\hbar F_1^{(1)}(y)\}$ is of order $\hbar^2$; at one loop $\hbar F_1^{(1)}$ enters linearly. $(x,y)\mapsto[\mathcal R^{(2)}(x,y)]$ is a two-cocycle of $\mathfrak{psl}(4|4)$ with values in $H^{-1}(Q+\{I,-\})$, and solvability through arity two, for some $\hbar F_1^{(1)}$, is the vanishing of its class in $H^2(\mathfrak{psl}(4|4);H^{-1}(Q+\{I,-\}))$.
\item (iv) Under the classical symmetry condition at the weights of the pairs, $H^{-1}(Q+\{I,-\})$ there is spanned by the classes $[S^{(1)}(z)]$, the adjoint module, and the uniform directions $\varphi(x)=t_xS^{(1)}(x)$ units change the class by
\begin{equation}
[\mathcal R^{(2)}(x,y)]\longmapsto[\mathcal R^{(2)}(x,y)]+\big(t_x+t_y-t_{[x,y]}\big)[S^{(1)}([x,y])]\ \text{units},
\label{eq:c12-tshift}
\end{equation}
where $t_{[x,y]}S^{(1)}([x,y])$ is $\sum_zt_zS^{(1)}(z)$ over the components $z$ of $[x,y]$ in the blocks of Theorem \ref{thm:c12-H}.
\end{itemize}
\end{theorem}

\begin{proof}
(i) Lemma \ref{lem:c12-wz} (iii). (ii) Definition \ref{def:c12-ob}; for $\varphi=(Q+\{I,-\})\psi$ the brackets are exact by \eqref{eq:c12-deriv}. (iii) The arity-three component of $s$ on a degree-$1$ cochain $G$ with $G^{(0)}=G^{(1)}=0$ is $(Q+\{I,-\})G^{(3)}$ plus a signed Chevalley--Eilenberg differential of $G^{(2)}$, the signs forced as in Lemma \ref{lem:c12-wz} (ii) by $s^2=0$ on $G=sF^{(2)}$; in $H^{-1}$ the first term vanishes, and $S^{(3)}=0$. (iv) $(-1)^{|x|}\{S^{(1)}(x),S^{(1)}(y)\}=-S^{(1)}([x,y])-(Q+\{I,-\})S^{(2)}(x,y)$ (Lemma \ref{lem:c2-components}), so the change in (ii) for $\varphi=tS^{(1)}$ units is $-t_{[x,y]}S^{(1)}([x,y])+(t_x+t_y)\big(S^{(1)}([x,y])+(Q+\{I,-\})S^{(2)}(x,y)\big)$ units.
\end{proof}

\begin{theorem}[What the arity-two equations fix]
\label{thm:c12-fix}
Assume the classical symmetry condition, the solvability of the arity-one equations, and write $[\mathcal R^{(2)}(x,y)]_{t=0}=[S^{(1)}(\varrho(x,y))]$ units with $\varrho(x,y)\in\mathfrak{psl}(4|4)$ of weight $\lambda_x+\lambda_y$, the class at $t_x=0$ for every $x$. The arity-two equations along all pairs are solvable exactly when
\begin{equation}
\varrho(x,y)=t\cdot[x,y]-[t\cdot x,y]-[x,t\cdot y]\qquad\text{for all }x,y,\qquad t\cdot z:=t_zz,
\label{eq:c12-rho}
\end{equation}
that is, when the cocycle $\varrho$ with values in the adjoint module is the coboundary of the diagonal one-cochain $t\cdot$. Pair by pair:

\begin{tabular}{l|l|l}
pair & bracket & equation\\
\hline
$(\Delta,P_b)$ & $-2P_b$ & $\varrho=2t_\Delta P_b$\\
$(\Delta,\xi)$ & $-\xi$ & $\varrho=t_\Delta\,\xi$\\
$(\xi_+,\xi_-)$ & $P_{\Gamma(\xi_+,\xi_-)|_V}$ & $\varrho=(t_P-t_{\xi_+}-t_{\xi_-})P_{\Gamma}$\\
$(\xi_+,\eta_+)$ & components along $\Delta$, $\mathfrak{so}(V)$, $\mathfrak{so}(W_{\mathbb R}^\perp)$ & $\varrho_z=(t_z-t_{\xi_+}-t_{\eta_+})\,z$\\
$(\eta_+,\eta_-)$ & $K_{\Gamma(\eta_+,\eta_-)|_V}$ & $\varrho=(t_K-t_{\eta_+}-t_{\eta_-})K_{\Gamma}$\\
$(K_b,P_{b'})$ & $(b,b')\Delta+M_{bb'}$ & $\varrho_\Delta=(t_\Delta-t_K-t_P)(b,b')\Delta$, $\varrho_M=(t_M-t_K-t_P)M_{bb'}$\\
$(M_\omega,P_b)$, $(M_\omega,\xi)$, $(R_r,\xi)$ & $[M_\omega,P_b]$, $[M_\omega,\xi]$, $[R_r,\xi]$ & $\varrho=-t_M[M_\omega,P_b]$, $-t_M[M_\omega,\xi]$, $-t_R[R_r,\xi]$\\
$(\Delta,\eta)$, $(\Delta,K_b)$ & $\eta$, $2K_b$ & $\varrho=-t_\Delta\eta$, $-2t_\Delta K_b$\\
$(P_b,\eta)$, $(K_b,\xi)$ & $\rho(b)\eta$, $\rho(b)\xi$ & $\varrho=(t_\xi-t_P-t_\eta)\rho(b)\eta$, $(t_\eta-t_K-t_\xi)\rho(b)\xi$\\
$(\xi_+,\xi'_+)$, $(\xi_+,\eta_-)$, $(\xi,P_b)$, $(\eta_+,\eta'_+)$, $(P_b,P_{b'})$ & $0$ & $\varrho=0$\\
\end{tabular}
\tabcaption{The arity-two equations of Theorem \ref{thm:c12-fix}, and likewise with the chiralities exchanged; in the rows $(P_b,\eta)$ and $(K_b,\xi)$ the image has the opposite chirality and $t_\xi$, $t_\eta$ are the coefficients of the image block.\label{tab:c12-fix}}

Consequently:
\begin{itemize}
\item (i) $t_\Delta$, $t_M$ and $t_R$ are fixed, $t_\Delta$ overdetermined by $(\Delta,P_b)$, $(\Delta,\xi_\pm)$ and the $\Delta$-components of $(\xi_\pm,\eta_\pm)$;
\item (ii) the solutions $(t_P,t_{\xi_\pm},t_{\eta_\pm},t_K)$ form an affine plane over the two diagonal derivations of $\mathfrak{psl}(4|4)$: the inner one, $t_z\mapsto t_z+c\lambda_z$, generated by the field redefinition $\phi^{(0)}=c\,S^{(1)}(\Delta)$ units, and the outer one, $t_z\mapsto t_z+aB_z$, the direction $G_B$ of Definition \ref{def:c8-B};
\item (iii) the remaining relations are consistency conditions, consequences of the cocycle condition of Theorem \ref{thm:c12-eq2} (iii);
\item (iv) in a translation-covariant scheme $t_P=t_\Delta$ and $t_{\eta_\pm}=t_{\xi_\mp}$ up to the outer direction, the inner direction being excluded since $\phi^{(0)}=c\,S^{(1)}(\Delta)$ units shifts $\hbar F_1^{(1)}(P_b)$ by $2c\,S^{(1)}(P_b)$ units and $\hbar F_1^{(1)}(\Delta)$ by $0$;
\item (v) on the commuting pairs $\varrho=0$ holds automatically, $\mathfrak{psl}(4|4)$ having no element of the required weight and type, and the computed classes there test the classical symmetry condition; on two legs the sector coefficients enter separately, the $[S^{(1)}(z)_d]$ being independent in $H^{-1}(Q)$, and the two-leg parts of the commuting-pair equations impose sector uniformity.
\end{itemize}
\end{theorem}

\begin{proof}
\eqref{eq:c12-tshift}: the equation along $(x,y)$ is $[S^{(1)}(\varrho(x,y))]+(t_x+t_y-t_{[x,y]})[S^{(1)}([x,y])]=0$ in units, which is \eqref{eq:c12-rho} componentwise. The brackets are those of Proposition \ref{prop:c2-special}, Remark \ref{rem:c2-dilatation} and Remark \ref{rem:c7-KDelta}; $[\xi_+,\eta_-]=0$ because $S_+\otimes S_-\otimes W\otimes W$ contains no element of $\mathfrak{so}(V)\oplus\mathfrak{so}(W_{\mathbb R}^\perp)\oplus\mathbb C\Delta$. On the six unknowns $t_P,t_{\xi_\pm},t_{\eta_\pm},t_K$ the equations from $(\xi_\pm,\eta_\pm)$, $(\xi_+,\xi_-)$, $(\eta_+,\eta_-)$, $(K_b,P_{b'})$ have rank four, the sum of the first four being the fifth up to the values of $\varrho$; the homogeneous solutions are $t_z=c\lambda_z$ and $t_z=aB_z$, whose diagonal maps are derivations ($\lambda_{[x,y]}=\lambda_x+\lambda_y$, and $B$ is a derivation). The inner direction: $(s\phi)^{(1)}(x)=\{\phi^{(0)},S^{(1)}(x)\}=c\{S^{(1)}(\Delta),S^{(1)}(x)\}\equiv c\lambda_xS^{(1)}(x)$ units modulo exact terms; the outer one: Proposition \ref{prop:c8-laws} (c).
\end{proof}

\begin{theorem}[The sector coefficients along $\Delta$ from the arity-two equation]
\label{thm:c12-theta}
Let $\hbar F_1^{(1)}(\Delta)$ and $\hbar F_1^{(1)}(P_b)$ be the solutions of Theorem \ref{thm:c12-sol}. On two legs the arity-two equation along $(\Delta,P_b)$ reads
\begin{equation}
Q\,\hbar F_1^{(2)}(\Delta,P_b)\big|_2+2\sum_{d\in\{\perp,\lambda,V\}}\theta^\Delta_d\,S^{(1)}(P_b)_d\ \text{units}=\hbar\,\mathrm{Ob}^{(2)}(\Delta,P_b)\big|_2+Q(\text{local}),
\label{eq:c12-theta}
\end{equation}
independently of the $\theta^{P_b}_d$. The classes $[S^{(1)}(P_b)_d]$ are independent in $H^{-1}(Q)$ on two legs; hence, with $o_d$ the coefficient of $[S^{(1)}(P_b)_d]$ in $[\hbar\,\mathrm{Ob}^{(2)}(\Delta,P_b)|_2]$, the equation is solvable on two legs exactly when this class lies in the span of the $[S^{(1)}(P_b)_d]$, and then $\theta^\Delta_d=\tfrac12o_d$ in units.
\end{theorem}

\begin{proof}
In the third equation of \eqref{eq:c12-eqs}, $[\Delta,P_b]=-2P_b$ and $S^{(2)}(\Delta,P_b)=0$. Sector parts in different sectors have no common legs, and within one sector $\{S^{(1)}(\Delta)_d,S^{(1)}(P_b)_d\}=v.S^{(1)}(P_b)_d=2S^{(1)}(P_b)_d$, the weight of $S^{(1)}(P_b)_d$ being $5$ (Lemma \ref{lem:c9-SDelta} (i)). The $\theta$-terms are therefore $2\theta^\Delta_dS^{(1)}(P_b)_d$ from $\{\hbar F_1^{(1)}(\Delta),S^{(1)}(P_b)\}$, $2\theta^{P_b}_dS^{(1)}(P_b)_d$ from $\{S^{(1)}(\Delta),\hbar F_1^{(1)}(P_b)\}$ and $-2\theta^{P_b}_dS^{(1)}(P_b)_d$ from $\hbar F_1^{(1)}(-2P_b)$, in units. The particular part of $\hbar F_1^{(1)}(\Delta)$ is a combination of field-counting functionals, which commute with $S^{(1)}(P_b)$, and exact terms, whose brackets are exact. Independence: $Q$ preserves the sectors $(A|_{W_{\mathbb R}^\perp},A^*|_{W_{\mathbb R}^\perp})$, $(\lambda,\lambda^*)$ and $(A|_V,A^*|_V)\oplus(c,c^*)$. In the first there is no functional of degree $-2$, $\int f(x)(A^*|_{W_{\mathbb R}^\perp},A^*|_{W_{\mathbb R}^\perp})=0$ identically, $A^*$ being odd and the fibre pairing symmetric. In the second the image is spanned by $Q\int(\rho(b)\lambda^*,\lambda^*)=2\int(\partial_b\lambda,\lambda^*)+2\int(\rho(b\wedge\partial)\lambda,\lambda^*)$, which is independent of $S^{(1)}(P_b)_\lambda$. In the third the translation-invariant image is spanned by $Q\int(\iota_bA|_V,c^*)$, which differs from $S^{(1)}(P_b)_V$ by the nonzero $\int(\iota_bdA|_V,A^*|_V)$ up to sign.
\end{proof}

\begin{theorem}[The computed arity-two equations on two legs]
\label{thm:c12-ar2}
With the solutions of Theorem \ref{thm:c12-sol}, the classes of the arity-two residuals on two legs vanish exactly under the conditions of Table \ref{tab:c12-ar2}. The obstructions are
\begin{equation}
\begin{aligned}
\hbar\,\mathrm{Ob}^{(2)}(\Delta,P_b)|_2&=\Big[-\tfrac23\int(\partial_bA|_{W_{\mathbb R}^\perp},A^*|_{W_{\mathbb R}^\perp})+\tfrac23\int(\partial_b\lambda,\lambda^*)+\big(-\tfrac32+\log\tfrac43\big)\int(\rho(b){\not d}\lambda,\lambda^*)\\
&\qquad-\tfrac23\int(\partial_bA|_V,A^*|_V)+\big(\tfrac32-\log\tfrac43\big)\int(\partial_\mu(\iota_bA|_V),A^*_\mu)+\big(-\tfrac56+\log\tfrac43\big)\int(\partial_bc,c^*)\Big]\ \text{units},\\
\hbar\,\mathrm{Ob}^{(2)}(\Delta,\xi)|_2&=\Big[8(T_1+T_2)-\tfrac{26}3(T_3+T_4)+\log\tfrac43\,Q(\psi_1+\psi_2)\Big]\ \text{units},\qquad \psi_1:=\int x^\mu\big(\Gamma(\xi,\lambda)^\mu,c^*\big),\quad \psi_2:=\int x^\mu\big(A^*_\mu,(\xi,\lambda^*)\big),
\end{aligned}
\label{eq:c12-ob2}
\end{equation}
with $o_d=-\tfrac23$ units in each sector, no position part and no pole along $(\Delta,P_b)$; $\hbar F_1^{(2)}(\Delta,\xi)|_2=\log\tfrac43\,(\psi_1+\psi_2)$ units plus closed directions of degree $-2$; $\hbar\,\mathrm{Ob}^{(2)}(\xi,P_b)|_2$ is $\tfrac14\log\tfrac43$ times a $Q$-exact functional; every two-leg arity-two pole sum vanishes. Consequently
\begin{equation}
t_\Delta=-\tfrac13,\qquad t_M=t_\Delta,\qquad t_R=t_\Delta+\tfrac13,\qquad t_{\xi_+}+t_{\xi_-}-t_P=\tfrac{27}4+2\log\tfrac32,\qquad t_{\xi_+}+t_{\eta_+}-t_\Delta=\tfrac{27}4+2\log\tfrac32,
\label{eq:c12-t}
\end{equation}
and the sector coefficients of $t_P$, $t_\xi$, $t_\eta$, $t_M$, $t_R$ are uniform. The value $t_\Delta=-\tfrac13$ is obtained twice, from $(\Delta,P_b)$ and from $(\Delta,\xi_\pm)$.
\end{theorem}

\begin{tabular}{l|l}
pair & condition\\
\hline
$(\Delta,P_b)$ & $\theta^\Delta_\perp=\theta^\Delta_\lambda=\theta^\Delta_V=-\tfrac13$\\
$(\Delta,\xi_\pm)$ & $\theta^\Delta_\perp=\theta^\Delta_\lambda=\theta^\Delta_V=-\tfrac13$, $\theta^\xi_d$ free\\
$(\xi_+,\xi_-)$ & $\theta^{P_b}_d=\theta^{\xi_+}_V+\theta^{\xi_-}_V-\tfrac{27}4-2\log\tfrac32$ for every $d$, and $\theta^{\xi_+}_V+\theta^{\xi_-}_V=\theta^{\xi_+}_\perp+\theta^{\xi_-}_\perp$\\
$(\xi_\pm,\xi'_\pm)$ & $\theta^{\xi_\pm}_V=\theta^{\xi_\pm}_\perp$\\
$(\xi_\pm,P_b)$ & $\theta^{P_b}_V=\theta^{P_b}_\perp=\theta^{P_b}_\lambda$\\
$(\xi_+,\eta_-)$ & $\theta^\xi_V-\theta^\xi_\perp=\theta^\eta_V-\theta^\eta_\perp$\\
$(\xi_+,\eta_+)$ & $\theta^R_\perp=\theta^R_\lambda=:\tau$, $\theta^\Delta_d=\theta^M_d=\tau-\tfrac13$, $\theta^\xi_d+\theta^\eta_d=\tau+\tfrac{77}{12}+2\log\tfrac32$ ($d=V,\perp$)\\
\end{tabular}
\tabcaption{The conditions under which the two-leg arity-two equations are solvable, in the parametrisation of \eqref{eq:c12-sol}.\label{tab:c12-ar2}}

\begin{proof}
The cycles of Proposition \ref{prop:c12-classes2} evaluated by \eqref{eq:c12-special}, the cohomology on complete bases fixed by representation counts, and the residuals of Definition \ref{def:c12-ob}; along $(\xi,\eta)$ solvability is decided on the value vectors. Theorem \ref{thm:c12-theta} gives the first row from $o_d=-\tfrac23$ units. With $\theta^\xi_V=\theta^\xi_\perp=:t_\xi$ the rows $(\xi_+,\eta_\pm)$ give $\theta^\eta_V=\theta^\eta_\perp=:t_\eta$, the position part of $\hbar F_1^{(1)}(\eta)$ then being $x^\mu\hbar F_1^{(1)}(\rho(e_\mu)\eta)$ with one coefficient; inserting $\theta^\Delta=-\tfrac13$ into the row $(\xi_+,\eta_+)$ gives $\tau=0$, and the rows give \eqref{eq:c12-t}.
\end{proof}

\begin{corollary}[Scheme-independent content on the action side]
\label{cor:c12-inv}
The inner direction ($t_z\mapsto t_z+c\lambda_z$, $\lambda_z=-1,+1,-2,+2$ on $\Sigma$, $\Sigma^*$, $P_b$, $K_b$ and $0$ on $\Delta$, $M_\omega$, $R_r$) and the outer direction ($t_{\xi_\pm}\mapsto t_{\xi_\pm}\pm a$, $t_{\eta_\pm}\mapsto t_{\eta_\pm}\mp a$) leave invariant exactly
\[
t_\Delta=-\tfrac13,\qquad t_M-t_\Delta=0,\qquad t_R-t_\Delta=\tfrac13,\qquad t_{\xi_+}+t_{\xi_-}-t_P=\tfrac{27}4+2\log\tfrac32,\qquad t_{\xi_+}+t_{\eta_+}-t_\Delta=\tfrac{27}4+2\log\tfrac32,
\]
together with the sector uniformity of $t_P$, $t_\xi$, $t_\eta$; these are the one-loop parts of $[D,P]=-2P$, $[D,Q]=-Q$, $[Q,P]=0$, $[Q,Q]=P$ and $\{Q,S\}=D+M+R$ in the form of Theorem \ref{thm:c12-algebra}. In a translation-covariant scheme $t_P=t_\Delta=-\tfrac13$, hence $t_{\xi_+}+t_{\xi_-}=\tfrac{77}{12}+2\log\tfrac32$ and $t_{\eta_+}=t_{\xi_-}$. The individual $t_{\xi_\pm}$, $t_{\eta_\pm}$, $t_P$, $t_K$, the coefficients $\theta^K_d$, $\kappa^K_j$, $\kappa_V$, $\kappa_\lambda$, the particular parts with $\log\tfrac32$ and every representative $\hbar F_1^{(k)}$ depend on the scheme.
\end{corollary}

\begin{proof}
Theorem \ref{thm:c12-ar2} and Theorem \ref{thm:c12-fix} (ii): each combination has zero weight under $\lambda$ and under $B$; for instance $t_{\xi_+}+t_{\xi_-}-t_P$ changes by $c(-1-1+2)$ and $a(1-1)$, and $t_{\xi_+}+t_{\eta_+}-t_\Delta$ by $c(-1+1)$ and $a(1-1)$. The translation-covariant scheme is Theorem \ref{thm:c12-fix} (iv).
\end{proof}

\begin{theorem}[One-loop algebra relations at the level of the action]
\label{thm:c12-algebra}
Assume the classical symmetry condition and the arity-one equations. For each pair $(x,y)$ the following are equivalent: (a) the arity-two equation along $(x,y)$ is solvable; (b) $[\mathcal R^{(2)}(x,y)]=0$, that is, \eqref{eq:c12-rho} holds at $(x,y)$; (c) with the one-loop corrected generators $S_\hbar(z):=S^{(1)}(z)+t_z\,S^{(1)}(z)$ units, $z\in\{x,y,[x,y]\}$, the part of order $\hbar$ of $(-1)^{|x|}\{S_\hbar(x),S_\hbar(y)\}+S_\hbar([x,y])$ plus $S^{(1)}(\varrho(x,y))$ units is $(Q+\{I,-\})$-exact. The relations are the one-loop corrections of $[Q,Q]=P$ along $(\xi,\xi')$, of $\{Q,S\}=D+M+R$ along $(\xi,\eta)$, of $[K,P]=D+M$ along $(K_b,P_{b'})$, of $[D,Q]=-Q$ along $(\Delta,\xi)$, of $[D,P]=-2P$ along $(\Delta,P_b)$ and of $\{S,S\}=K$ along $(\eta,\eta')$. Computed on two legs: $[D,P]$, $[D,Q]$, $[Q,P]$, $[Q,Q]$ and $\{Q,S\}$, by Theorem \ref{thm:c12-ar2}.
\end{theorem}

\begin{proof}
(a)$\Leftrightarrow$(b) is Theorem \ref{thm:c12-eq2} (i). (b)$\Leftrightarrow$(c): with $(-1)^{|x|}\{S^{(1)}(x),S^{(1)}(y)\}=-S^{(1)}([x,y])-(Q+\{I,-\})S^{(2)}(x,y)$, the part of order $\hbar$ is $(t_x+t_y-t_{[x,y]})S^{(1)}([x,y])$ units modulo exact terms, and \eqref{eq:c12-tshift} applies.
\end{proof}

\begin{remark}[The coefficient $\theta$ as a cross-check]
\label{rem:c12-thetacheck}
Theorem \ref{thm:c9-kappa} derives $\theta=-\tfrac13$ from the arity-two component of $L\circ L=0$ on the scalar bilinears, under the hypothesis that the quantisation and the local representatives exist through arity two, and Corollary \ref{cor:c11-theta} with Proposition \ref{prop:c11-theta-sector} does the same family by family for the three sector coefficients. In the present construction $\theta^\Delta_d=-\tfrac13$ is computed from the arity-two obstruction along $(\Delta,P_b)$ and again along $(\Delta,\xi_\pm)$ (Theorem \ref{thm:c12-ar2}), and agrees with these derivations.
\end{remark}

